\newpage
\begin{center}
{\fontsize{16pt}{19.2pt}\selectfont \textbf{BUILDING A REINFORCEMENT LEARNING MODEL FOR OPTIMIZING REORDER POLICY IN MULTI-WAREHOUSE INVENTORY MANAGEMENT}\par}
\vspace{0.5cm}
{\fontsize{16pt}{19.2pt}\selectfont \textbf{ABSTRACT}\par}
\end{center}
\vspace{0.2cm}

Inventory control across a retail network of multiple warehouses and multiple
stock-keeping units (SKUs) is a sequential decision-making problem under
uncertainty that requires balancing several conflicting cost components at
once: holding cost, stockout cost, fixed ordering cost, warehouse overflow
penalty, and the requirement to sustain a target customer service level.
Classical policies such as EOQ, $(s, S)$ and the Newsvendor model rest on
restrictive assumptions --- deterministic or stationary demand, fixed lead
times, and unconstrained storage capacity --- and therefore lose their
optimality guarantees once demand is stochastic, lead times fluctuate, and
multiple SKUs must share a finite warehouse capacity.

This thesis addresses the problem through multi-agent reinforcement learning.
A system of $N = 10$ warehouses and $M = 50$ SKUs is modelled as $500$ agents,
one per warehouse--SKU pair, each issuing its own replenishment decisions from
local observations. The agents are trained with Independent Multi-Agent
Proximal Policy Optimization (IPPO) under full parameter sharing: a single
Actor--Critic network is shared across all $500$ agents, which keeps the number
of learnable parameters constant as the system scales. The reward is designed
with local decomposition and a 30-day rolling-window fill-rate penalty term,
reducing the variance of the learning signal relative to a single shared
global reward.

The simulation environment follows the Gymnasium interface, with demand
parameters calibrated from the real-world M5 Walmart retail dataset (10
stores, 50 SKUs, 1,941 days of sales data) and stochastic lead times drawn
from one to three days. The proposed model is benchmarked against three
classical inventory policies (EOQ, $(s, S)$, Newsvendor) within the same
environment under a unified evaluation protocol involving multiple random
seeds and statistical significance testing, across four metric groups: total
operational cost, fill rate, and the breakdown of individual cost components.

Experimental results are analysed along three dimensions: (i) the comparison
between the IPPO agent and the classical baselines in terms of total
operational cost and service level; (ii) the generalization capability of the
shared-parameter policy across SKUs of differing demand scale; and (iii) the
effect of local reward decomposition on the learning process. The thesis also
reports scenarios in which the proposed method does not outperform the
classical baselines, together with an analysis of the underlying causes.

